Let \(P\in\mathcal P(\alpha,\beta,s,d,M_Y,L,c,t_0)\), and put \(m=\lfloor\alpha\rfloor\). By membership in the class of Definition~\texttt{def:holder-class}, for every \(x\in[0,1]^d\) the function \(a\mapsto\mu_P(a,x)\) lies in the same univariate Hölder ball of order \(\alpha\) and radius \(L\), uniformly in \(x\). In particular, under the stated Hölder convention, each derivative \(\partial_a^j\mu_P(a,x)\) required for \(0\le j\le m\) is jointly measurable and uniformly bounded in absolute value by a constant \(C_{\mathrm H}L\), where \(C_{\mathrm H}\) depends only on the convention.

We first transfer the derivatives through the probability integral. For \(j=1,\ldots,m\), suppose inductively that
\[
\theta_P^{(j-1)}(a)=\int\partial_a^{j-1}\mu_P(a,x)p_X(x)\,dx.
\]
For every sufficiently small nonzero \(r\), the mean-value theorem and the uniform derivative bound give
\[
\left|\frac{\partial_a^{j-1}\mu_P(a+r,x)-\partial_a^{j-1}\mu_P(a,x)}{r}\right|\le C_{\mathrm H}L.
\]
The difference quotient converges pointwise in \(x\) to \(\partial_a^j\mu_P(a,x)\). Since \(p_X(x)\,dx\) is a probability measure, dominated convergence yields
\[
\theta_P^{(j)}(a)=\int\partial_a^j\mu_P(a,x)p_X(x)\,dx,
\qquad
|\theta_P^{(j)}(a)|\le C_{\mathrm H}L.
\]
At an endpoint, the identical argument uses the corresponding one-sided difference quotients. Beginning with \(\theta_P(a)=\int\mu_P(a,x)p_X(x)\,dx\), induction proves the display for every derivative required by the convention. Continuity of these derivatives follows as well: if \(a_r\to a\), pointwise continuity of \(\partial_a^j\mu_P(\cdot,x)\), together with the same uniform bound, permits another application of dominated convergence.

If \(\alpha\notin\mathbb N\), write \(\gamma=\alpha-m\in(0,1)\). The uniform Hölder remainder bound and integration against the probability density give, for all \(a,a'\in[0,1]\),
\[
\begin{aligned}
|\theta_P^{(m)}(a)-\theta_P^{(m)}(a')|
&\le \int |\partial_a^m\mu_P(a,x)-\partial_a^m\mu_P(a',x)|p_X(x)\,dx\\
&\le C_{\mathrm H}L|a-a'|^\gamma.
\end{aligned}
\]
If \(\alpha\in\mathbb N\), the bounded-derivative requirement in the frozen convention has already been transferred by the preceding display. Thus every seminorm entering the univariate Hölder norm of \(\theta_P\) is bounded by \(C_{\mathrm H}L\). Consequently \(a\mapsto\theta_P(a)\) belongs to a univariate Hölder ball of order \(\alpha\), with radius depending only on \(L\) and the chosen Hölder-norm convention. No boundedness, smoothness, or positivity condition on \(p_X\) was used beyond \(p_X(x)\,dx\) being the probability law \(P_X\).